\section{1. Introduction}

Light field (LF) imaging captures both spatial and angular information of light rays, offering significant advantages in 3D reconstruction, autonomous navigation, and robotic vision \cite{lightfieldimaging3Dreconst}. Accurate depth estimation from LF data constitutes a primary task in these applications, enabling precise scene understanding and spatial interaction \cite{depthestimationsceneunderst}. By exploiting the rich angular samples, conventional LF depth estimation methods achieve high accuracy in controlled environments.

However, most current depth estimation frameworks rely strictly on the Lambertian assumption, which presumes angular photo-consistency across different viewpoints \cite{Lambertianassumptionphotoco}. This assumption holds for matte surfaces but fails when encountering non-Lambertian materials, such as specular or scattering surfaces, which are ubiquitous in real-world environments \cite{nonLambertianmaterialsspecu}. Unfortunately, light reflection over these complex surfaces violates the working principles of standard Epipolar Plane Image (EPI) based algorithms \cite{EpipolarPlaneImageEPI}. Specifically, specular highlights shift inconsistently across views, destroying the linear EPI structures required for reliable disparity calculation.

The severity of this limitation becomes evident through quantitative evaluations of standard EPI-based depth estimation networks. Our baseline experiments demonstrate that while the model achieves a low Mean Absolute Error (MAE) of 0.081 in mixed urban scenes, its performance degrades substantially in purely non-Lambertian scenes, yielding an unacceptable MAE of 0.411. Visual analysis confirms that the network predictions in non-Lambertian regions completely collapse into mere replications of the input texture, entirely missing the underlying geometric depth. Such severe performance degradation renders existing LF depth estimation techniques unreliable for practical deployment in complex scenarios.

Addressing this challenge remains difficult due to the intrinsic complexity of mixed pixels and the severe scarcity of annotated non-Lambertian LF datasets. When multiple reflection types coexist within a single pixel, standard hard-classification strategies fail to capture the continuous material transitions. Furthermore, purely data-driven approaches suffer from severe overfitting under limited non-Lambertian samples, with standard convolutional classifiers achieving a mere 24.6\% accuracy in material identification. Therefore, it is highly necessary to develop a unified, physics-driven framework that transcends the Lambertian constraint, enabling robust depth estimation across diverse reflective environments.

Conventional light field depth estimation methods primarily exploit epipolar plane images or focal stacks to extract geometric structures \cite{EPIfocalstacks}. These approaches strictly rely on the Lambertian assumption, which presumes angular photo-consistency across different viewpoints \cite{Lambertianassumption}. However, this assumption fails when encountering non-Lambertian materials, such as specular or scattering surfaces \cite{nonLambertianmaterials}. Specifically, specular highlights shift inconsistently across views, destroying the linear structures required for reliable disparity calculation \cite{specularhighlights}. To mitigate this, some heuristic methods attempt to detect and mask specular pixels using clustering techniques \cite{heuristicmasking}. Unfortunately, these handcrafted rules struggle in mixed scenes where multiple materials coexist at the sub-pixel level. Furthermore, they cannot model scattering materials, inevitably producing depth maps with severe artifacts and missing regions \cite{scatteringmaterials}.

Early deep learning-based methods treat light fields as multi-view stacks, employing convolutional neural networks to implicitly map features to depth \cite{earlydeeplearningEPINet}. While achieving high accuracy on synthetic Lambertian datasets, these data-driven models lack explicit physical optics priors \cite{physicalopticspriors}. Consequently, these models suffer from severe overfitting when trained on limited non-Lambertian samples with scarce annotations. For instance, standard learning-based material classifiers often yield extremely low accuracy when processing small angular patches. Technically, the fixed receptive fields of standard convolutions cannot adapt to non-linear photometric variations across views. As a result, the network predictions in non-Lambertian regions completely collapse into mere replications of input textures. This texture replication phenomenon significantly degrades the depth estimation accuracy, yielding unacceptably high errors in complex scenes.

Recent advanced approaches attempt to address these limitations by introducing attention mechanisms or explicit reflection models \cite{attentionmechanismsreflectio}. Some methods decouple the light field into diffuse and specular layers using separate network branches \cite{decouplelightfield}. However, these techniques typically rely on hard classification or independent processing pipelines for different materials. This discrete treatment ignores the physical reality of mixed pixels, where continuous material transitions frequently occur. Moreover, incorporating complex rendering losses or heavy pre-trained backbone networks substantially increases the computational burden. Such computationally expensive designs significantly exceed the $O(HW)$ processing limits of resource-constrained or real-time applications. Furthermore, these complex architectures exhibit poor generalization when facing domain shifts or unseen materials in real environments. Ultimately, they fail to provide a unified theoretical framework that simultaneously explains diffuse, non-Lambertian, and mixed scenes.

In summary, existing methods either lack physical interpretability, struggle with small-sample material classification, or fail to resolve texture replication. Therefore, there is an urgent need for a unified, physically grounded approach that explicitly models angular frequency distributions. Such a method must dynamically adapt depth estimation strategies based on continuous physical priors. This ensures robust performance across diverse and complex real-world scenes without relying on massive data-driven training.

In this paper, we propose a unified dual-mask physical model for non-Lambertian light field depth estimation. Instead of relying solely on data-driven feature extraction, our approach explicitly integrates physical optics priors into the deep learning framework. Specifically, we introduce a dual-mask physical decomposition module that predicts diffuse and specular masks to decouple non-Lambertian reflections. These masks subsequently guide the construction of an epipolar cost volume, effectively suppressing erroneous matches in highlight regions. Finally, a 3D regularization network processes the physically consistent cost volume to regress continuous and accurate depth maps. Our main contributions are summarized as follows:

\begin{itemize}\item \textbf{Contribution 1}: We propose an MRI-inspired angular frequency analysis mechanism that transforms spatial angular distributions into the frequency domain via 2D discrete Fourier transform. This mechanism classifies materials based on spectral energy distributions, distinguishing diffuse, specular, and scattering reflections without requiring massive training data. It resolves the overfitting bottleneck of learning-based classifiers on small non-Lambertian samples, providing reliable physical priors with $\mathcal{O}(HW)$ computational complexity.\end{itemize}

\begin{itemize}\item \textbf{Contribution 2}: We develop a unified dual-mask physical modeling framework coupled with an inverse wavevector parsing strategy. By designing a medium mask for roughness quantification and an angular direction mask for wavevector variation, we establish a comprehensive physical optics modulation model. This framework outputs continuous material weights rather than hard classifications, effectively addressing the coexistence of multiple materials in mixed pixels and enhancing model interpretability.\end{itemize}

\begin{itemize}\item \textbf{Contribution 3}: We introduce a component-aware adaptive depth estimation strategy that dynamically selects estimation algorithms based on the parsed material weights. We construct a mask-guided epipolar cost volume and incorporate pixel-level material weights into the loss function to suppress non-Lambertian noise. This strategy effectively resolves the texture replication failure of conventional epipolar plane image methods in non-Lambertian regions, significantly improving depth accuracy in complex mixed scenes.\end{itemize}

Extensive experiments on four benchmark datasets, including HCInew and UrbanLF-Syn, demonstrate the superiority of our method. Our model achieves an overall mean absolute error of 0.133 on mixed datasets, strictly meeting the target thresholds for both Lambertian and non-Lambertian sub-datasets. These results validate that integrating explicit physical models with deep networks yields robust and accurate depth estimation for real-world complex environments.
